\addcontentsline{toc}{chapter}{СПИСОК ИСПОЛЬЗОВАННЫХ ИСТОЧНИКОВ}
\begin{thebibliography}{}
	\bibitem{python3}  Python /  [Электронный ресурс] // Python : [сайт]. — URL: https://www.python.org/ (дата обращения: 24.09.2024).
	\bibitem{python3-time}  time — Time access and conversions /  [Электронный ресурс] // Python 3.12.6 documentation : [сайт]. — URL: https://docs.python.org/3/library/time.html\#time.process\_time\_ns (дата обращения: 24.09.2024).
	\bibitem{python3-matplotlib}  Matplotlib, a library for creating visualizations in Python /  [Электронный ресурс] // Matplotlib : [сайт]. — URL: https://https://matplotlib.org/ (дата обращения: 25.09.2024).
	\bibitem{tracemalloc}  tracemalloc — Trace memory allocations /  [Электронный ресурс] // Python 3.12.6 documentation : [сайт]. — URL: https://docs.python.org/3/library/tracemalloc.html\#tracemalloc.get\_traced\_memory (дата обращения: 20.09.2024).
        \bibitem{algs} Томас Х. Кормен, Чарльз И. Лейзерсон, Рональд Л. Ривест, Клиффорд Штайн Алгоритмы: построение и анализ — 3-е изд. — М.: «Вильямс», 2013. — с. 440. — ISBN 978-5-8459-1794-2
        \bibitem{lev} В. И. Левенштейн (1965). Двоичные коды с исправлением выпадений, вставок и замещений символов. Доклады Академии Наук СССР. 163 (4): 845–848.
\end{thebibliography}